\section{Industriesteuerung / Controller}

Die Industriesteuerung stellt die zentrale Einheit des Systems dar. 
Sie verarbeitet sämtliche eingehenden Signale von den verschiedenen Plattformen (Tablet, Laptop, VR) sowie den physischen Schaltern und steuert darauf basierend die Ausgänge, wie beispielsweise die Beleuchtung im Modellhaus. 
Die Kommunikation zwischen den Systemkomponenten erfolgt über eine OPC-UA-Schnittstelle.

\subsection{Hardware-Konfiguration}

Die Steuerung wird mit einer Versorgungsspannung von \textbf{24V DC} betrieben. 
Die Spannungsversorgung erfolgt über eine digitale Ein- und Ausgangskarte.

Die Eingänge werden zur Erfassung von Schaltsignalen verwendet, während die Ausgänge zur Ansteuerung der Verbraucher (z.B. Glühbirne) dienen.

% Bild hier einfügen
% \includegraphics[width=0.6\textwidth]{dein_bild.png}
% \caption{Digitale Ein- und Ausgangskarte}

\subsection{IO-Mapping}

Das IO-Mapping definiert die Zuordnung zwischen physikalischen Ein- und Ausgängen und deren Funktion im System.

\begin{table}[h]
\centering
\begin{tabular}{|c|c|c|}
\hline
\textbf{Signal} & \textbf{Adresse} & \textbf{Beschreibung} \\
\hline
Schalter Steckbrett & I0.0 & Manueller Lichtschalter \\
Tablet & I0.1 & Steuerung über Tablet \\
Laptop & I0.2 & Steuerung über PC-Anwendung \\
VR-Headset & I0.3 & Steuerung über VR \\
\hline
Lampe & Q0.0 & Ausgang zur Glühbirne \\
\hline
\end{tabular}
\caption{IO-Mapping der Steuerung}
\end{table}

\subsection{Ethernet-Konfiguration}

Die Industriesteuerung ist über Ethernet mit dem Netzwerk verbunden. 
Dadurch wird die Kommunikation mit externen Anwendungen ermöglicht.

Die IP-Adresse der Steuerung wird statisch konfiguriert, sodass eine stabile Verbindung zu den Clients gewährleistet ist. 
Über das Netzwerk können Befehle empfangen sowie Zustände zurückgemeldet werden.

\subsection{OPC UA Konfiguration}

Für die Kommunikation zwischen Steuerung und den verschiedenen Plattformen wird das \textbf{OPC UA Protokoll} verwendet.

OPC UA ermöglicht:
\begin{itemize}
    \item Plattformunabhängige Kommunikation
    \item Echtzeit-Datenübertragung
    \item Synchronisierung zwischen allen Systemen
\end{itemize}

Die Steuerung fungiert hierbei als Server, während die Anwendungen (Tablet, Laptop, VR) als Clients agieren.

\subsection{User-Konfiguration}

Für den Zugriff auf die Steuerung werden Benutzer angelegt. 
Diese ermöglichen eine kontrollierte Verbindung über OPC UA.

Typische Einstellungen:
\begin{itemize}
    \item Benutzername und Passwort
    \item Zugriffsrechte auf Variablen
\end{itemize}

\subsection{OPC UA Daten pro Raum}

Für jeden Raum wird eine eigene Datenstruktur definiert, um die Zustände der Lampen zu verwalten.

\begin{table}[h]
\centering
\begin{tabular}{|c|c|}
\hline
\textbf{Variable} & \textbf{Bedeutung} \\
\hline
Licht\_Kueche & Status der Lampe (0 = aus, 1 = ein) \\
Licht\_Wohnzimmer & Status der Lampe \\
Licht\_Schlafzimmer & Status der Lampe \\
\hline
\end{tabular}
\caption{OPC UA Variablen pro Raum}
\end{table}

\subsection{Programm pro Raum}

Die Logik der Steuerung basiert auf einer ODER-Verknüpfung der Eingänge. 
Das bedeutet, dass die Lampe eingeschaltet wird, sobald eines der Eingangssignale aktiv ist.

\subsubsection{Logikbeschreibung}

Die Lampe wird aktiviert, wenn mindestens eine der folgenden Bedingungen erfüllt ist:
\begin{itemize}
    \item Schalter betätigt
    \item Signal von Tablet
    \item Signal von Laptop
    \item Signal vom VR-Headset
\end{itemize}

\subsubsection{Beispielhafte Implementierung}

\begin{verbatim}
IF (Schalter OR Tablet OR Laptop OR VR) THEN
    Lampe := TRUE;
ELSE
    Lampe := FALSE;
END_IF
\end{verbatim}

Zusätzlich wird eine positive Flankenerkennung verwendet, um nur auf Zustandsänderungen zu reagieren und unerwünschte Mehrfachschaltungen zu vermeiden.

\subsection{Zusammenfassung}

Die Industriesteuerung übernimmt die zentrale Verarbeitung aller Eingangssignale und sorgt für eine konsistente Ansteuerung der Ausgänge. 
Durch die Verwendung von OPC UA wird eine stabile und plattformübergreifende Kommunikation zwischen realem Modell und digitalen Anwendungen gewährleistet.